\begin{tabularx}{\textwidth}{L{23mm}X}
\toprule
\textbf{Jelölés} & \textbf{Jelentés} \\
\midrule
$\Sigma$ & Az alaphalmaz, amelynek elemein a műveletet végezzük. \\
$\oplus$ & Bináris, asszociatív művelet. Prefixhez és redukcióhoz az asszociativitás a kulcsfeltétel. \\
$e$ & A művelet semleges eleme, ha létezik. Például összeadásnál $0$, szorzásnál $1$. \\
$X=[x_1,\ldots,x_n]$ & A bemeneti sorozat. \\
$Y=[y_1,\ldots,y_n]$ & A prefixszámítás kimeneti sorozata. \\
$p$ & A számítási egységek száma. \\
$T_p(n)$ & Párhuzamos futási idő $p$ számítási egység mellett. \\
$W(n)$ & Teljes munka, vagyis az összes elvégzett elemi művelet száma. \\
\bottomrule
\end{tabularx}
